\documentclass[11pt]{amsart}

\usepackage[utf8]{inputenc}
\usepackage{amsmath,amssymb,amsthm}
\usepackage[hidelinks]{hyperref}

\theoremstyle{plain}
\newtheorem{theorem}{Theorem}
\newtheorem{lemma}[theorem]{Lemma}
\newtheorem{proposition}[theorem]{Proposition}
\newtheorem{conjecture}[theorem]{Conjecture}
\newtheorem{claim}[theorem]{Claim}
\newtheorem{question}[theorem]{Question}

\def\bE{\mathbb{E}}
\def\bR{\mathbb{R}}
\def\Pb{\mathbb{P}}

\begin{document}

\title{On the Probability a Weighted Bernoulli Sum Exceeds Its Mean}
\author{}
\date{}
\maketitle

\section{Statement}

Let $w\in\bR_{>0}^m$ be a vector of positive weights with $\sum_{i=1}^m w_i=1$, and
let $v\in\{0,1\}^m$ be a random binary vector whose coordinates are i.i.d.\
$\mathrm{Bernoulli}(p)$. Set $X=\langle w,v\rangle$. Then $X$ takes values in $[0,1]$
and $\bE[X]=p$. We study the following question.

\begin{question}\label{q:prob}
For which $p\in(0,1)$ does the following hold? For every $w\in\bR_{>0}^m$ with
$\sum_{i=1}^m w_i=1$, if $v\in\{0,1\}^m$ has i.i.d.\ $\mathrm{Bernoulli}(p)$
coordinates, then
\[
\Pb[\langle v,w\rangle\ge p]\ge p .
\]
\end{question}

We relate Question~\ref{q:prob} to the Manickam--Mikl\'os--Singhi (MMS) conjecture.

\begin{conjecture}[MMS]\label{conj:MMS}
Let $n,k$ be positive integers with $n\ge 4k$. For any real numbers
$x_1,\dots,x_n$ with $\sum_{i=1}^n x_i=0$, at least $\binom{n-1}{k-1}$ of the
$k$-element subsets $S\subseteq[n]$ satisfy $\sum_{i\in S}x_i\ge 0$.
\end{conjecture}

A linear bound (the conjecture holds for all $n\ge ck$ with an absolute constant
$c$, e.g.\ $c=10^{46}$) is known by Pokrovskiy. The aim of this note is the
following conditional theorem.

\begin{theorem}\label{thm:main}
Suppose there is a constant $c>0$ such that Conjecture~\ref{conj:MMS} holds for
all positive integers $n,k$ with $n\ge ck$. Then Question~\ref{q:prob} holds for
all $p\in[0,1/c)$.
\end{theorem}

In what follows we only ever apply Conjecture~\ref{conj:MMS} through its hypotheses
$\sum_i x_i=0$ and $n\ge ck$; no structural restriction on the values $x_i$ is
used.

\section{Proof of Theorem~\ref{thm:main}}

The case $p=0$ is trivial (the inequality $\Pb[X\ge 0]\ge 0$ always holds), so
assume $p>0$. We first prove the statement for rational $p\in(0,1/c)$, and then
extend it to all real $p\in(0,1/c)$ by a limiting argument.

\subsection{Two elementary facts}

We isolate two standard facts used below.

\begin{lemma}[Left-continuity of the upper tail]\label{lem:leftcont}
Let $Z$ be a real-valued random variable and define
$g(s)=\Pb[Z\ge s]$ for $s\in\bR$. Then $g$ is left-continuous: if
$s_1\le s_2\le\cdots$ and $s_i\to s$, then $g(s)=\lim_{i\to\infty}g(s_i)$.
\end{lemma}

\begin{proof}
Since $s_i\le s$, the event $\{Z\ge s\}$ is contained in $\{Z\ge s_i\}$ for every
$i$, and the events $\{Z\ge s_i\}$ are nested decreasing in $i$ (as $s_i$ is
nondecreasing). We claim
\[
\{Z\ge s\}=\bigcap_{i\ge1}\{Z\ge s_i\}.
\]
The inclusion ``$\subseteq$'' is immediate. For ``$\supseteq$'', if $Z\ge s_i$ for
all $i$, then $Z\ge\sup_i s_i=\lim_i s_i=s$. By continuity of probability from
above for a decreasing sequence of events,
$g(s)=\Pb\big[\bigcap_i\{Z\ge s_i\}\big]=\lim_i\Pb[Z\ge s_i]=\lim_i g(s_i)$.
\end{proof}

\begin{lemma}[Monotone coupling]\label{lem:coupling}
Fix $w\in\bR_{>0}^m$ and $0<q\le p<1$. Let $v$ have i.i.d.\
$\mathrm{Bernoulli}(p)$ coordinates and let $v'$ have i.i.d.\
$\mathrm{Bernoulli}(q)$ coordinates. Then for every threshold $s\in\bR$,
\[
\Pb[\langle w,v\rangle\ge s]\ \ge\ \Pb[\langle w,v'\rangle\ge s].
\]
\end{lemma}

\begin{proof}
Let $U_1,\dots,U_m$ be i.i.d.\ uniform on $[0,1]$ and set
$v_j=\mathbf 1[U_j\le p]$ and $v'_j=\mathbf 1[U_j\le q]$. Then $v$ has i.i.d.\
$\mathrm{Bernoulli}(p)$ coordinates, $v'$ has i.i.d.\
$\mathrm{Bernoulli}(q)$ coordinates, and since $q\le p$ we have $v'_j\le v_j$ for
every $j$ pointwise. As all $w_j>0$,
$\langle w,v\rangle\ge\langle w,v'\rangle$ pointwise on this probability space.
Hence $\{\langle w,v'\rangle\ge s\}\subseteq\{\langle w,v\rangle\ge s\}$, which
gives the claimed inequality.
\end{proof}

\subsection{The rational case}

Fix a rational $p=k_0/n_0\in(0,1/c)$ with positive integers $k_0,n_0$, and fix a
weight vector $w\in\bR_{>0}^m$ with $\sum_{j=1}^m w_j=1$. Let $v$ have i.i.d.\
$\mathrm{Bernoulli}(p)$ coordinates and $X=\langle w,v\rangle$. We must show
$\Pb[X\ge p]\ge p$.

\paragraph{A combinatorial model.}
Set $k=k_0$ and $n=n_0$, so that $k/n=p$. Construct a random $0/1$ matrix $M$ with
$n$ rows and $m$ columns as follows: independently for each column
$j\in[m]$, place a single $1$ in a uniformly random row and $0$'s elsewhere.
Denote the rows of $M$ by $r_1,\dots,r_n\in\{0,1\}^m$ and put
$y_i=\langle w,r_i\rangle$ for $i\in[n]$. Independently of $M$, choose a uniformly
random $k$-element subset $S\subseteq[n]$, and set $Y=\sum_{i\in S}y_i$.

\begin{claim}\label{cl:dist}
$Y$ has the same distribution as $X=\langle w,v\rangle$.
\end{claim}

\begin{proof}
We have $Y=\big\langle w,\ \sum_{i\in S}r_i\big\rangle$, so it suffices to show that
$u:=\sum_{i\in S}r_i$ has the same distribution as $v$. Condition on a fixed
$k$-set $S$. For each column $j\in[m]$, the $j$-th coordinate of $u$ equals $1$
iff the unique row containing the $1$ of column $j$ lies in $S$; since that row is
uniform over $[n]$ and the columns are placed independently, the coordinates of
$u$ are independent and each equals $1$ with probability $|S|/n=k/n=p$. Thus,
conditionally on $S$, $u$ has i.i.d.\ $\mathrm{Bernoulli}(p)$ coordinates,
identical to the law of $v$. As this conditional law does not depend on $S$, it is
also the unconditional law of $u$. Hence $u\sim v$ and $Y\sim X$.
\end{proof}

\begin{claim}\label{cl:bound}
$\Pb[Y\ge p]\ge p$.
\end{claim}

\begin{proof}
Define $z_i=y_i-\frac1n$ for $i\in[n]$. Since each column of $M$ contains exactly
one $1$, we have $\sum_{i=1}^n r_i=\mathbf 1$ (the all-ones vector in
$\bR^m$), hence
\[
\sum_{i=1}^n y_i=\Big\langle w,\sum_{i=1}^n r_i\Big\rangle
=\langle w,\mathbf 1\rangle=\sum_{j=1}^m w_j=1,
\]
and therefore
\begin{equation}\label{eq:sumzero}
\sum_{i=1}^n z_i=\sum_{i=1}^n y_i-n\cdot\frac1n=0 .
\end{equation}
Note that \eqref{eq:sumzero} holds for \emph{every} realization of $M$.

Since $p=k/n<1/c$, we have $n>ck$, hence $n\ge ck$. Condition on any realization
of $M$; this fixes the real numbers $z_1,\dots,z_n$, which satisfy
\eqref{eq:sumzero}. Applying the hypothesis (Conjecture~\ref{conj:MMS} for
$n\ge ck$) to $z_1,\dots,z_n$, at least $\binom{n-1}{k-1}$ of the $k$-element
subsets $S\subseteq[n]$ satisfy $\sum_{i\in S}z_i\ge0$. Because $S$ is chosen
uniformly among all $\binom{n}{k}$ such subsets, independently of $M$,
\[
\Pb\Big[\textstyle\sum_{i\in S}z_i\ge0 \ \Big|\ M\Big]
\ \ge\ \frac{\binom{n-1}{k-1}}{\binom{n}{k}}
=\frac{k}{n}=p ,
\]
and this bound holds for every realization of $M$. Taking expectation over $M$,
\begin{equation}\label{eq:avg}
\Pb\Big[\textstyle\sum_{i\in S}z_i\ge0\Big]\ \ge\ p .
\end{equation}
Finally, $\sum_{i\in S}z_i=\sum_{i\in S}y_i-\frac{|S|}{n}=Y-\frac{k}{n}=Y-p$, so
the event $\{\sum_{i\in S}z_i\ge0\}$ coincides with $\{Y\ge p\}$. Substituting
into \eqref{eq:avg} gives $\Pb[Y\ge p]\ge p$.
\end{proof}

Combining Claims~\ref{cl:dist} and~\ref{cl:bound},
\[
\Pb[X\ge p]=\Pb[Y\ge p]\ge p ,
\]
which proves Question~\ref{q:prob} for every rational $p\in(0,1/c)$.

\subsection{Extension to irrational $p$}

Let $p\in(0,1/c)$ be arbitrary (in particular possibly irrational), fix
$w\in\bR_{>0}^m$ with $\sum_j w_j=1$, let $v$ have i.i.d.\
$\mathrm{Bernoulli}(p)$ coordinates, and set $X=\langle w,v\rangle$. Choose a
sequence of rationals $p_1\le p_2\le\cdots$ with $0<p_i<p$ and $p_i\to p$; this is
possible since $p>0$.

For each $i$, the rational case applied with parameter $p_i$ (and the same $w$)
gives: if $v^{(i)}$ has i.i.d.\ $\mathrm{Bernoulli}(p_i)$ coordinates, then
\begin{equation}\label{eq:rat}
\Pb[\langle w,v^{(i)}\rangle\ge p_i]\ \ge\ p_i .
\end{equation}
Since $p_i\le p$, Lemma~\ref{lem:coupling} (with $q=p_i$ and threshold $s=p_i$)
yields
\begin{equation}\label{eq:dom}
\Pb[\langle w,v\rangle\ge p_i]\ \ge\ \Pb[\langle w,v^{(i)}\rangle\ge p_i]\ \ge\ p_i ,
\end{equation}
where the last inequality is \eqref{eq:rat}. Thus $\Pb[X\ge p_i]\ge p_i$ for all
$i$.

Now apply Lemma~\ref{lem:leftcont} to the random variable $Z=X$ and the increasing
sequence $s_i=p_i\to p$:
\[
\Pb[X\ge p]=\lim_{i\to\infty}\Pb[X\ge p_i]\ \ge\ \lim_{i\to\infty}p_i=p ,
\]
where the inequality uses \eqref{eq:dom}. This proves Question~\ref{q:prob} for
$p$.

Since $p\in(0,1/c)$ was arbitrary and the case $p=0$ is trivial, the theorem
follows. \qed

\section{Optimality outside the threshold}

For completeness we record that the answer to Question~\ref{q:prob} is negative
for $p\in(1/3,1)\setminus\{1/2\}$, so a positive answer can hold only on an initial
segment together with the isolated point $p=1/2$.

\begin{proposition}\label{prop:counterexample}
The answer to Question~\ref{q:prob} is negative for every $p\in(1/3,1)$ with
$p\ne1/2$.
\end{proposition}

\begin{proof}
Suppose first $1/3<p<1/2$ and take $w=(1/3,1/3,1/3)$. Then
$\langle w,v\rangle\ge p$ iff at least two coordinates of $v$ equal $1$ (since the
attainable values of $\langle w,v\rangle$ are $0,1/3,2/3,1$ and $1/3<p<2/3$).
Hence
\[
\Pb[\langle w,v\rangle\ge p]=\binom{3}{2}p^2(1-p)+p^3=3p^2-2p^3=p(3p-2p^2).
\]
Then
\[
p-\Pb[\langle w,v\rangle\ge p]=p\big(1-3p+2p^2\big)=p(1-2p)(1-p)>0
\]
for $p\in(1/3,1/2)$, so $\Pb[\langle w,v\rangle\ge p]<p$.

If $1/2<p<1$, take $w=(1/2,1/2)$. Then $\langle w,v\rangle\ge p$ iff both
coordinates of $v$ equal $1$ (the attainable values are $0,1/2,1$ and $1/2<p<1$),
so $\Pb[\langle w,v\rangle\ge p]=p^2<p$.
\end{proof}

\end{document}
